% Source: https://artofproblemsolving.com/community/c3h2685536_slopeintercept_with_circles
% Stable identity: aops-topic-2685536-q01
The circles do meet, since $CD=\sqrt{4^2+2^2}=2\sqrt5<12$.

\emph{Method 1 (subtract the equations).} The circles are $(x-5)^2+(y-4)^2=36$ and $(x-1)^2+(y-2)^2=36$. Subtracting,
\[\left[(x-5)^2-(x-1)^2\right]+\left[(y-4)^2-(y-2)^2\right]=(-8x+24)+(-4y+12)=0.\]
Both $A$ and $B$ satisfy both circle equations, hence this linear equation, so it is the line $AB$: $\boxed{y=-2x+9}$.

\emph{Method 2 (perpendicular bisector).} $A$ and $B$ are each at distance $6$ from both $C$ and $D$, so both lie on the perpendicular bisector of $\overline{CD}$. The midpoint of $\overline{CD}$ is $M(3,3)$ and the slope of $CD$ is $\dfrac{4-2}{5-1}=\dfrac12$, so the perpendicular slope is $-2$. The line is $y-3=-2(x-3)$, i.e. $y=-2x+9$ again.

\begin{asy}
usepackage("amsmath");
size(6cm);
defaultpen(fontsize(10pt));
pair C=(5,4),D=(1,2),M=(3,3);
pair u=(1,-2)/sqrt(5);pair A=M+sqrt(31)*u,B=M-sqrt(31)*u;
draw((-5.5,0)--(11.5,0),gray,Arrow(4));draw((0,-4.5)--(0,10.5),gray,Arrow(4));
draw(circle(C,6));draw(circle(D,6));
draw((-0.6,10.2)--(6.6,-4.2),linewidth(0.9));
draw(C--D,dashed);dot(C);dot(D);dot(M);dot(A);dot(B);
pair v1=unit(C-M),v2=unit(A-M);draw(M+0.6*v1--M+0.6*(v1+v2)--M+0.6*v2);
label("$C(5,4)$",C,NE,fontsize(8pt));label("$D(1,2)$",D,W,fontsize(8pt));label("$M$",M,dir(160),fontsize(8pt));
label("$A$",A,E);label("$B$",B,NE);label("$y=-2x+9$",(6.3,-3.6),E,fontsize(8pt));
\end{asy}
